\documentclass[aspectratio=169,10pt]{beamer}

\usetheme{NMA}

\title[Models]{Models are useful abstractions}
\subtitle{A compact tour of descriptive, mechanistic, and normative views}
\author[Gunnar Blohm]{Gunnar Blohm}
\date{Neuromatch Academy}

%\course{Computational Neuroscience}
%\week{Week 1}
%\session{Day 1}
%\topic{Project}

\begin{document}

\begin{frame}[plain]
  \titlepage
\end{frame}

\begin{frame}{Why models are useful}
  \begin{itemize}
    \item Synthesize knowledge and expose hidden assumptions.
    \item Retrieve latent information that cannot be measured directly.
    \item Turn hypotheses into quantitative predictions.
    \item Guide experiments, interventions, and new technologies.
  \end{itemize}

  \NMAtakeaway{A model is useful when it sharpens the question we can ask.}
  \end{frame}

\begin{frame}{Match the abstraction to the question}
  \begin{columns}[T,onlytextwidth]
    \begin{column}{.47\textwidth}
      \begin{block}{Descriptive -- ``what?''}
        Summarize patterns in data with the smallest useful representation.
      \end{block}
      \vspace{.35em}
      \begin{exampleblock}{Mechanistic -- ``how?''}
        Specify a process that can generate the observed behavior.
      \end{exampleblock}
    \end{column}
    \begin{column}{.47\textwidth}
      \begin{alertblock}{Normative -- ``why?''}
        Ask what an ideal system should do under explicit constraints.
      \end{alertblock}
      \vspace{.35em}
      Each view answers a different question; no single level is always best.
    \end{column}
  \end{columns}
\end{frame}

\begin{frame}{Cosine tuning captures a broad pattern}
  \begin{columns}[c,onlytextwidth]
    \begin{column}{.42\textwidth}
      A cosine tuning model relates firing rate to movement direction:
      \[
        r(\theta)=b+A\cos(\theta-\theta_0).
      \]
      \begin{itemize}
        \item $b$: baseline activity
        \item $A$: modulation depth
        \item $\theta_0$: preferred direction
      \end{itemize}
    \end{column}
    \begin{column}{.53\textwidth}
      \centering
      \begin{tikzpicture}[x=1cm,y=1cm,>=stealth]
        \draw[->,NMAMidGray] (0,0) -- (5.1,0) node[right] {$\theta$};
        \draw[->,NMAMidGray] (0,0) -- (0,3.0) node[above] {$r$};
        \draw[NMATeal,very thick,smooth,samples=80,domain=0:4.8]
          plot (\x,{1.45+1.05*cos((\x-1.15)*75)});
        \draw[dashed,NMARule] (1.15,0) -- (1.15,2.5);
        \node[anchor=south,text=NMACharcoal] at (1.15,2.52) {preferred direction};
      \end{tikzpicture}
    \end{column}
  \end{columns}
\end{frame}

\NMAsectionframe[Models and data form a cycle of discovery]{From description to mechanism}

\begin{frame}{Models and experiments constrain each other}
  \begin{columns}[c,onlytextwidth]
    \begin{column}{.36\textwidth}
      \begin{itemize}
        \item Models turn ideas into predictions.
        \item Experiments test those predictions.
        \item New data revise the model and the question.
      \end{itemize}
    \end{column}
    \begin{column}{.58\textwidth}
      \centering
      \begin{tikzpicture}[
        node distance=1.55cm and 1.70cm,
        every node/.style={font=\small,align=center},
        box/.style={rounded corners=2pt,fill=NMATeal,text=white,
                    minimum width=2.15cm,minimum height=.80cm},
        arrow/.style={->,thick,NMAMidGray,rounded corners=5pt}
      ]
        \node[box] (question) {Question\\Hypothesis};
        \node[box,right=of question] (model) {Model\\Quantification};
        \node[box,below=of model] (prediction) {Prediction\\Outcome};
        \node[box,left=of prediction] (experiment) {Experiment\\New data};
        \draw[arrow] (question) -- (model);
        \draw[arrow] (model) -- (prediction);
        \draw[arrow] (prediction) -- (experiment);
        \draw[arrow] (experiment) -- (question);
      \end{tikzpicture}
    \end{column}
  \end{columns}
\end{frame}

\begin{frame}{Use figures as evidence, not decoration}
  \begin{columns}[T,onlytextwidth]
    \begin{column}{.47\textwidth}
      \begin{block}{Preferred asset formats}
        \begin{itemize}
          \item PDF for plots and diagrams
          \item SVG converted to PDF at build time
          \item TikZ for simple native schematics
        \end{itemize}
      \end{block}
    \end{column}
    \begin{column}{.47\textwidth}
      \begin{alertblock}{Reserve raster images for}
        Photographs, microscopy, or source material that is inherently pixel-based.
      \end{alertblock}
    \end{column}
  \end{columns}

  \vfill
  {\small\color{NMAMidGray}%
    Keep source and attribution near the figure, where the audience can read it.}
\end{frame}

\begin{frame}{Design choices should support learning}
  \begin{enumerate}
    \item Start with the question the audience should be able to answer.
    \item Choose the simplest model that makes the relevant distinction.
    \item Show the evidence and state what it changes.
  \end{enumerate}

  \vspace{.5em}
  \begin{block}{Template principle}
    One claim per slide, generous margins, and a consistent visual hierarchy.
  \end{block}
\end{frame}

\NMAclosingframe[Questions and productive next steps]{Now build a model worth testing}

\end{document}
